\documentclass{syllabus}

\begin{document}

\thispagestyle{fancy}
\coursecode{PHY 2200}
\renewcommand{\lms}{Brightspace}
\date{Fall 2026}
\author{}
\title{\coursecode: General Physics I \\ {\large \mdseries Introduction to Mechanics}}
\maketitle

\meeting{MWF: & 1:35-3:55 PM, DSC 095}
\officehours{MWF: & 2:00-3:30 PM}
% TODO: confirm the 5/e year. 4/e was Wiley 2015; 5/e added Titus and
% Spicklemire as authors.
\textbook{Matter and Interactions, Vol. I (5/e)}{Chabay and Sherwood}{Wiley}{2025}
\prerequisites{MTH 1120: & (Calculus I) \\ & MTH 1220: & (Calculus II, co-enrolled)}


\section*{General Information}

\infotable
\hspace{\fill}
\includegraphics[height=2.25in]{Assets/Chabay-cover.png}
\hspace{.25in}

\section*{Course overview}

This course is an introduction to mechanics. What we mean by this is that it introduces you to a particular way of studying the relationship between forces on physical bodies and the motion of those bodies. In the sense that everyone has deduced some principles of how forces and motion are related, everyone is a physicist. However, the particular way of studying these principles which you will learn in this class was expressed by Isaac Newton over three hundred years ago as follows:

\begin{quote}
    \itshape Mechanics can be considered in a twofold respect: as rational, which proceeds \textbf{accurately by demonstration}; and practical. To practical mechanics all the \textbf{manual arts} belong, from which mechanics took its name\ldots{} In this sense rational mechanics will be the science of motions resulting from any forces whatsoever, and of the forces required to produce any motions, accurately proposed and demonstrated\ldots{} And therefore we offer this work as the \textbf{mathematical principles of philosophy}; for all the difficulty of philosophy seems to consist in this---from the phenomena of motions to investigate the forces of nature and then from these forces to demonstrate the other phenomena.
    \upshape
    \begin{flushright}
        Isaac Newton, in \textit{The Mathematical Principles of Natural Philosophy} (1687)
    \end{flushright}
\end{quote}

You already know practical mechanics. The purpose of this class is threefold: to identify, formalize, and extend the physics you \textit{already know} through practical means. We will distill principles in precise mathematical language, and we will see how those principles can be extended to predict surprising results about forces and motion.

\subsection*{Student learning outcomes}

The following concepts are important for organizing the course materials and your own understanding of mechanics:
\begin{enumerate}[label=\textcolor{CarthageDark}{\bfseries\arabic*.}]
    \item Use the concepts of body and quantity to develop principles of interaction between objects.
    \item Differentiate between various descriptions of motion and the situations in which they are applicable.
    \item Appraise the effectiveness of approximations to physical laws and identify possible forms of higher-order corrections to physical laws.
    \item Apply fundamental physics principles, laws, and theories to predict real-world phenomena and evaluate the reasonableness of these predictions.
    \item Set up and solve physics problems both analytically (with symbols) and computationally.
    \item Measure key quantities associated with physical phenomena and use physics principles to interpret the resulting data.
\end{enumerate}

\section*{Course components}

\paragraph{Homework} Homework offers critically important practice for new physics ideas and methods, and helps form the foundation upon which deeper physics understanding depends. Homework for this course is administered and graded through WileyPLUS, reached from the course {\lms} page, so that you get immediate feedback on your work. There is no separate access code to purchase or class key to enter.

\paragraph{Extensions} Automatic extensions on homework can be requested up to 10 days beyond the due date and are active for 24 hours after they are requested. A \textit{20\% late penalty will be applied to points earned after the assignment's due date has passed.} Extensions beyond the 10-day window must be manually enabled and are subject to instructor approval. To request a manual extension for an assignment, send the instructor an email with an explanation for the need and the new requested deadline.

\paragraph{Labs} Labs are directed mini-experiments (lab studios) that arise naturally from newly learned material. These exercises typically involve physics measurements, inquiries, calculations, and analyses in teams of 2 or 3 students. Given the frequency of lab studios, the equipment they use, and their group nature, \textit{some labs cannot be easily made up if missed}. Monitor the class schedule to make sure you do not miss lab days.

\paragraph{VPython} Computer programming is an essential skill for physics of the 21st century. Very few interesting problems can be solved with pen and paper, and even calculators provide little additional leverage. For this reason, \coursecode{} incorporates programming and simple computer modeling using the Glowscript/VPython environment. Glowscript is a relatively easy-to-learn computer programming language designed for 3D visualizations and physics-capable calculations. \coursecode{} is not a class in computer science, however. You will not be graded on the elegance of your codes but rather on their functionality, the correctness of the physics, and whether the code fulfills the assignment.

\paragraph{Quizzes} Short quizzes will be administered approximately once per week at the beginning of the class period. Grading criteria emphasize the procedures used to derive solutions, not the circled answers at the end. \textit{All work must be shown for full credit to be awarded.} After the first quiz or two, key equations will be provided on a summary sheet. \textit{You may use a calculator and one sheet of handwritten notes during quizzes and exams.} No other electronic devices are permitted unless instructed otherwise for that quiz.

\paragraph{Exams} Two cumulative mid-term exams will be administered during the semester. Their format will be largely similar to quizzes but longer and cumulative in topical coverage. A comprehensive final exam will take place during the final exam period.

\section*{Grading Scheme}
\subsection*{Component Weights}

\noindent\begin{tabular}{ll}
\textbf{Component} & \textbf{Weight}       \\ \hline
\textbf{Homework}               & 20\%  \\
\textbf{Unit quizzes}           & 24\%  \\
\textbf{Lab studios}            & 10\%  \\
\textbf{VPython}                & 10\%  \\
\textbf{Exams}                  & 36\%
\end{tabular}%

% TODO: confirm these criteria. The F24 syllabus left this sentence unfinished,
% so this clause is borrowed from PHY 1200.
\lettergrades{a) a pattern of seeking help in the course at office hours, b) number of missing homework assignments.}

\section*{Policies}
\subsection{Electronic equipment}
Electronic equipment can be a distraction to you and others during in-person instruction. Please be considerate and restrict use of phones and non-class related use of tablets and laptops to break periods. Unless otherwise instructed, phones, tablets, and laptops may not be used during assessments.

\subsection{General course policies}
Expect a response within one business day for emails sent to the instructor. To ensure a prompt response, format the subject line as: [\coursecode] *insert subject here*

Students are expected to arrive and leave on time. I will attempt to arrive five minutes early and be available for at least five minutes after class.

\collegepolicies

\section*{Course resources}

\paragraph{Course website} The class site, where lecture notes and study resources are available for inspection and download.

\paragraph{WileyPLUS} The homework and gradebook system for this course, reached from the course {\lms} page. Access to the etext is included.

\paragraph{Course schedule} The detailed course schedule is a separate document that shows day-to-day class activities, reading assignments, lecture topics, exam dates, and quiz dates. Please check this document regularly to keep up to date.

\paragraph{Web VPython (glowscript.org)} VPython is a Python computing language plus a 3D graphics module used for in-class numerical simulations. Glowscript is a cloud computing resource that saves your work automatically and enables VPython programming without installing anything on your computer.

\tipsforsuccess{upper}

\section*{Course Schedule}

\noindent\begin{tabular}{cll}
\textbf{Unit} & \textbf{Topic} & \textbf{Weeks} \\ \hline
0 & Vectors                 & 0     \\
1 & Forces and motion       & 1-3   \\
2 & Energy and total motion & 3-8   \\
3 & Curved motion           & 8-10  \\
4 & Physics of matter       & 11-13 \\
5 & Atomic physics          & 13-14
\end{tabular}%

\end{document}
